\documentclass[11pt,a4paper]{article}
\usepackage[margin=3cm]{geometry}
\usepackage{amsmath,amssymb,amsthm,mathrsfs}
\usepackage[T1]{fontenc}
\usepackage{microtype}
\usepackage{booktabs}
\usepackage[colorlinks=true,linkcolor=blue,citecolor=blue]{hyperref}

\theoremstyle{plain}
\newtheorem{theorem}{Theorem}[section]
\newtheorem{proposition}[theorem]{Proposition}
\newtheorem{lemma}[theorem]{Lemma}
\newtheorem{corollary}[theorem]{Corollary}
\theoremstyle{definition}
\newtheorem{definition}[theorem]{Definition}
\newtheorem{remark}[theorem]{Remark}
\newtheorem{hypothesis}[theorem]{Hypothesis}

\newcommand{\R}{\mathbb{R}}
\newcommand{\E}{\mathbb{E}}
\newcommand{\Sw}{\mathcal S'}
\newcommand{\ip}[2]{\langle #1,#2\rangle}
\DeclareMathOperator{\Var}{Var}
\DeclareMathOperator{\Ent}{Ent}
\DeclareMathOperator{\supp}{supp}

\title{(T3): marginalisation of the log-Sobolev inequality to sharp time\\
\large and which Dirichlet form the programme actually needs}
\author{}
\date{}

\begin{document}
\maketitle

\begin{abstract}
(T3) is settled. If the space-cutoff Euclidean measure $\nu_g$ on $\Sw(\R^2)$ satisfies a Poincar\'e
or log-Sobolev inequality with respect to the \emph{covariance-weighted} Dirichlet form
$\E_\nu\ip{\nabla F}{G\nabla F}$, then its sharp-time marginal $\mu_g$ satisfies the same
inequality, \emph{with the same constant}, with respect to $\E_\mu\ip{\nabla f}{C\nabla f}$,
$C=(2\mu)^{-1}$ (Theorem~\ref{thm:T3}). Two exact facts drive it: the sharp-time restriction of
$G=(-\Delta_2+m_0^2)^{-1}$ \emph{is} $C$ (Lemma~\ref{lem:sharptime}), and time-mollification only
\emph{decreases} the Dirichlet form, so the inequality survives the limit in the favourable
direction (Lemma~\ref{lem:moll}). Both verified numerically (\S\ref{sec:num}).

\smallskip
\noindent\textbf{A sharper question for (T1).} The main theorem of \texttt{h6-attack.tex} requires
the \emph{covariance-weighted} form. The plain $L^2$ Dirichlet form is \textbf{useless} for it: for
a Wick monomial of degree $n\ge2$, $\E\bigl[|\nabla F_n[u]|_{L^2}^2\bigr]=+\infty$ identically,
because $\E[({:}\varphi^m(x){:})^2]=m!\,C(0)^m=\infty$ (Proposition~\ref{prop:L2fails}).
Covariance-weighted LSI implies $L^2$-LSI but not conversely. So (T1) must verify not merely
\emph{that} Bauerschmidt--Dagallier prove an LSI, but \emph{in which form}. That is now the
sharpest single question in the programme.
\end{abstract}

\section*{Update (8 August 2026)}

The conditional marginalisation theorem below remains a statement about a
\emph{covariance-weighted} two-dimensional inequality.  The question posed in the abstract has
now been answered by checking \cite[Eq.~(1.6) and Thm.~1.1]{BD}: their LSI uses the unweighted
lattice $\ell^2$ gradient.  It therefore does not satisfy hypothesis (P-2d) and does not feed
Theorem~\ref{thm:T3}.  No comparison reverses this mismatch because the Fourier multiplier of
$C=(2\sqrt{-\partial_x^2+m_0^2})^{-1}$ tends to zero at high momentum.  Thus (T3) is a valid
conditional lemma, but the proposed BD $\Rightarrow$ (P-2d) $\Rightarrow$ (P-U) route is closed.
See \texttt{remaining.tex} for the corrected program.

\section{Statement}

Let $G:=(-\Delta_{\R^2}+m_0^2)^{-1}$, $\Phi_G$ the centred Gaussian measure on $\Sw(\R^2)$ with
covariance $G$, and $d\nu_g=Z_g^{-1}e^{-U_g}d\Phi_G$ with
$U_g=\lambda\int_{\R^2}{:}P(\varphi(t,x)){:}\,g(x)\,dt\,dx$ the space-cutoff Euclidean measure; by
Feynman--Kac--Nelson its sharp-time marginal under $\pi:\varphi\mapsto\varphi(0,\cdot)$ is the
ground-state measure $\mu_g=\pi_*\nu_g=\Omega_g^2\,d\varphi_0$. Write $\mu=(-\partial_x^2+m_0^2)^{1/2}$
and $C=(2\mu)^{-1}$.

\begin{definition}[covariance-weighted Dirichlet form]\label{def:form}
For a centred Gaussian reference with covariance $\mathcal C$ and $\nu\ll\Phi_{\mathcal C}$, put
$\mathcal E_{\mathcal C}^\nu(F):=\E_\nu\ip{\nabla F}{\mathcal C\,\nabla F}$, $\nabla$ the
Cameron--Martin gradient. This is the form for which the Gaussian itself satisfies Poincar\'e with
constant one, $\Var_{\Phi_{\mathcal C}}(F)\le\E_{\Phi_{\mathcal C}}\ip{\nabla F}{\mathcal C\nabla F}$.
\end{definition}

\begin{hypothesis}[(P-2d)]\label{hyp:P2d}
There is $c_P<\infty$, independent of $g$, such that, for all $F$ in the form domain,
\[
 \Var_{\nu_g}(F)\le c_P\,\E_{\nu_g}\ip{\nabla F}{G\nabla F}.
\]
\end{hypothesis}

\section{The sharp-time covariance identity}

\begin{lemma}\label{lem:sharptime}
For every $k\in\R$, with $\mu(k)=(k^2+m_0^2)^{1/2}$,
\[
  \int_\R\frac{dk_0}{2\pi}\ \frac{1}{k_0^2+k^2+m_0^2}\;=\;\frac{1}{2\mu(k)} .
\]
Consequently $\pi_*\Phi_G=\varphi_0$, the centred Gaussian on $\Sw(\R)$ with covariance
$C=(2\mu)^{-1}$; equivalently $G\bigl((0,x),(0,y)\bigr)=C(x-y)$.
\end{lemma}

\begin{proof}
$\int_\R dk_0\,(k_0^2+a^2)^{-1}=\pi/a$ with $a=\mu(k)>0$; divide by $2\pi$. A centred Gaussian is
determined by its covariance, and the covariance of $\pi_*\Phi_G$ has Fourier symbol
$\int\frac{dk_0}{2\pi}(k_0^2+\mu(k)^2)^{-1}$.
\end{proof}

\section{Marginalisation}

Fix a smooth probability density $\rho\ge0$ on $\R$ with $\int\rho=1$, and set
$\rho_\varepsilon(t)=\varepsilon^{-1}\rho(t/\varepsilon)$ and
$\varphi_\varepsilon(x):=\int\rho_\varepsilon(t)\varphi(t,x)\,dt$. For a bounded smooth cylinder
functional $f$ on $\Sw(\R)$ put $F_\varepsilon(\varphi):=f(\varphi_\varepsilon)$.

\begin{lemma}[mollification decreases the form]\label{lem:moll}
$\displaystyle \ip{\nabla F_\varepsilon}{G\,\nabla F_\varepsilon}
 =\int\frac{dk}{2\pi}\,|\widehat{v_\varepsilon}(k)|^2\,s_\varepsilon(k)$, where
$v_\varepsilon:=(\nabla f)(\varphi_\varepsilon)$ and
\[
  s_\varepsilon(k):=\int\frac{dk_0}{2\pi}\,\frac{|\hat\rho(\varepsilon k_0)|^2}{k_0^2+\mu(k)^2} .
\]
Moreover $0\le s_\varepsilon(k)\le\widehat C(k)=\tfrac1{2\mu(k)}$ for every $\varepsilon>0$, and
$s_\varepsilon(k)\to\widehat C(k)$ as $\varepsilon\downarrow0$. Hence
\begin{equation}\label{eq:dominate}
  \ip{\nabla F_\varepsilon}{G\nabla F_\varepsilon}\ \le\ \ip{v_\varepsilon}{C\,v_\varepsilon}
  \qquad\text{pointwise, for every }\varepsilon>0 .
\end{equation}
\end{lemma}

\begin{proof}
$F_\varepsilon$ depends on $\varphi$ only through $\varphi_\varepsilon$, so
$\delta F_\varepsilon/\delta\varphi(t,x)=\rho_\varepsilon(t)\,v_\varepsilon(x)$, i.e.
$\nabla F_\varepsilon=\rho_\varepsilon\otimes v_\varepsilon$. Pairing with $G$ in Fourier and using
$\widehat{\rho_\varepsilon}(k_0)=\hat\rho(\varepsilon k_0)$ gives the displayed formula. Since
$\rho\ge0$ and $\int\rho=1$ we have $|\hat\rho|\le1$, whence
$s_\varepsilon(k)\le\int\frac{dk_0}{2\pi}(k_0^2+\mu^2)^{-1}=\widehat C(k)$ by
Lemma~\ref{lem:sharptime}; convergence is dominated convergence with the same integrable majorant,
using $\hat\rho(0)=1$ and continuity of $\hat\rho$.
\end{proof}

\begin{theorem}[(T3)]\label{thm:T3}
Assume \textup{(P-2d)}. Then for every bounded smooth cylinder functional $f$ on $\Sw(\R)$,
\[
  \Var_{\mu_g}(f)\ \le\ c_P\ \E_{\mu_g}\ip{\nabla f}{C\,\nabla f} ,
\]
with the \emph{same} constant $c_P$, uniformly in $g$. The same statement holds with $\Var$
replaced by $\Ent(\,\cdot^2)$ and $c_P$ by the log-Sobolev constant.
\end{theorem}

\begin{proof}
Apply (P-2d) to $F_\varepsilon$ and use \eqref{eq:dominate}:
\[
  \Var_{\nu_g}(F_\varepsilon)\ \le\ c_P\,\E_{\nu_g}\ip{\nabla F_\varepsilon}{G\nabla F_\varepsilon}
  \ \le\ c_P\,\E_{\nu_g}\ip{v_\varepsilon}{C\,v_\varepsilon} .
\]
The sharp-time restriction exists for $P(\varphi)_2$, so $\varphi_\varepsilon\to\varphi(0,\cdot)$ in
$\nu_g$-probability; as $f$ and $\nabla f$ are bounded and continuous, bounded convergence gives
$\Var_{\nu_g}(F_\varepsilon)\to\Var_{\nu_g}(f\circ\pi)=\Var_{\mu_g}(f)$ and
$\E_{\nu_g}\ip{v_\varepsilon}{Cv_\varepsilon}\to\E_{\mu_g}\ip{\nabla f}{C\nabla f}$. Passing to the
limit gives the claim. The entropy version is identical, the only change being that
$\Ent_{\nu_g}(F_\varepsilon^2)\to\Ent_{\mu_g}(f^2)$, which holds for $f$ bounded above and below
away from zero; the general case follows by the usual truncation.
\end{proof}

\begin{remark}[why the direction is favourable]
The whole content is \eqref{eq:dominate}: smearing in time can only \emph{shrink} the Dirichlet
form, because $|\hat\rho|\le1$. Had the inequality gone the other way, the limit
$\varepsilon\downarrow0$ would have destroyed the estimate. No constant is lost.
\end{remark}

\section{Which Dirichlet form: a sharper question for (T1)}\label{sec:which}

\begin{proposition}\label{prop:L2fails}
Let $F_n[u]=\int u\,{:}\varphi^n{:}$ with $u\not\equiv0$ and $n\ge2$. Then the \emph{unweighted}
Dirichlet form is identically infinite:
\[
  \E_{\varphi_0}\Bigl[\,\bigl|\nabla F_n[u]\bigr|_{L^2}^2\Bigr]
  =n^2\!\int u(x)^2\,\E\bigl[({:}\varphi(x)^{n-1}{:})^2\bigr]dx
  =n^2\,(n-1)!\,C(0)^{\,n-1}\!\int u^2\;=\;+\infty ,
\]
since $C(0)=+\infty$ logarithmically in one space dimension.
\end{proposition}

\begin{proof}
$\nabla F_n[u](x)=n\,u(x){:}\varphi(x)^{n-1}{:}$, and
$\E[({:}\varphi(x)^{m}{:})^2]=m!\,C(x,x)^m$ by the Wick isometry.
\end{proof}

\begin{remark}[consequence]\label{rem:consequence}
The main theorem of \texttt{h6-attack.tex} is driven by the identity
$\ip{\nabla F_n[u]}{C\nabla F_n[u]}
=n^2\iint u(x)u(y)C(x-y){:}\varphi(x)^{n-1}{:}{:}\varphi(y)^{n-1}{:}$: the factor $C(x-y)$
supplied by the \emph{weighting} is exactly what removes the coincident-point singularity that
Proposition~\ref{prop:L2fails} exhibits. Therefore:
\begin{itemize}
\item a covariance-weighted LSI \textbf{is} what the programme needs;
\item it implies the $L^2$ version, since
      $\ip{\nabla F}{C\nabla F}\le\|C\|_{\rm op}|\nabla F|^2_{L^2}=(2m_0)^{-1}|\nabla F|^2_{L^2}$;
\item the converse fails, because $\widehat C(k)=(2\mu(k))^{-1}\to0$ as $|k|\to\infty$, so $C$ is
      not bounded below and no reverse comparison exists.
\end{itemize}
The source check now shows that Bauerschmidt--Dagallier use the
$L^2$/stochastic-quantisation form, so their theorem cannot enter Theorem~\ref{thm:T3}.  Moreover,
even a covariance-weighted estimate would not by itself validate the former Theorem~2.3 of
\texttt{h6-attack.tex}; that proof has an independent divergent-trace error.  The corrected
conditional recursion and its additional hypotheses are in \texttt{remaining.tex},
Proposition~6.2.
\end{remark}

\section{Numerical verification}\label{sec:num}

Both analytic inputs checked numerically ($m_0=1$; $k_0$-integration on $[-4000,4000]$ with
$4\cdot10^6$ points; Gaussian mollifier $|\hat\rho(\varepsilon k_0)|^2=e^{-(\varepsilon k_0)^2}$).

\begin{center}
\begin{tabular}{rrrr}
\multicolumn{4}{c}{\textbf{(a) Lemma~\ref{lem:sharptime}}}\\
\toprule
$k$ & numeric & $1/(2\mu(k))$ & rel.\ error\\
\midrule
$0$   & $0.4999204$ & $0.5000000$ & $1.6\cdot10^{-4}$\\
$1$   & $0.3534738$ & $0.3535534$ & $2.2\cdot10^{-4}$\\
$10$  & $0.0496723$ & $0.0497519$ & $1.6\cdot10^{-3}$\\
\bottomrule
\end{tabular}
\qquad
\begin{tabular}{rrr}
\multicolumn{3}{c}{\textbf{(b) Lemma~\ref{lem:moll}}, $k=1$}\\
\toprule
$\varepsilon$ & $s_\varepsilon(1)$ & $\le\widehat C(1)$?\\
\midrule
$1.00$ & $0.118866$ & yes\\
$0.30$ & $0.232172$ & yes\\
$0.10$ & $0.303518$ & yes\\
$0.03$ & $0.337244$ & yes\\
$0.01$ & $0.347981$ & yes\\
\midrule
$0$ & $0.353553$ & ---\\
\bottomrule
\end{tabular}
\end{center}

The residual error in (a) is the truncation of the $1/k_0^2$ tail at $|k_0|=4000$, consistent in
size with $2\mu(k)/(\pi\cdot4000)$. In (b), $s_\varepsilon$ increases monotonically to
$\widehat C$ from below at every $k$ tested, as Lemma~\ref{lem:moll} asserts.

\section{Status}

\begin{itemize}
\item \textbf{(T3) is proved} (Theorem~\ref{thm:T3}), with no loss of constant, given (P-2d).
\item \textbf{The BD premise fails}: \cite{BD} proves the unweighted form, not (P-2d).  Hence its
      inequality cannot be passed through Theorem~\ref{thm:T3} to obtain (P-U).
\item Tasks (T1), (T2), and (T4) are irrelevant to this route unless a genuinely
      covariance-weighted two-dimensional estimate is first established.
\end{itemize}

\begin{thebibliography}{9}
\bibitem{BD} R.~Bauerschmidt and B.~Dagallier, \emph{Log-Sobolev inequality for the $\varphi^4_2$
and $\varphi^4_3$ measures}, Comm.\ Pure Appl.\ Math.\ \textbf{77} (2024) 2579--2612;
arXiv:2202.02295.
\bibitem{M3} Milestone M3, \texttt{../m3/m3.tex}.
\end{thebibliography}

\end{document}
